\chapter{Implementation}\label{ch:impl}

The authors' repository is used unedited, as a git submodule at \fpath{third_party/Table-Specialist/} (commit
\texttt{af1921f}); \texttt{git status} in the submodule is clean. Everything that runs is our code: 6,437 lines under
\fpath{ts/}, \fpath{scripts/} and \fpath{sandbox/}, plus 23 test files with 205 tests passing and 1 skipped. Repository
code enters our pipeline in only three ways (\cref{sec:provenance}), each through one module, \fpath{ts/repo.py}, or a
syntax-tree reader.

\section{Architecture}

\begin{figure}[htbp]
\centering
\begin{adjustbox}{max width=\linewidth}\begin{tikzpicture}[
  m/.style={draw=sea,fill=oursbg,rounded corners=2pt,thick,align=center,font=\sffamily\scriptsize,minimum height=0.8cm,text width=2.5cm},
  r/.style={m,draw=amber,fill=paperbg},
  lane/.style={draw=headshade,fill=rowshade!45,rounded corners=4pt,inner xsep=6pt,inner ysep=6pt},
  lab/.style={font=\sffamily\scriptsize\bfseries,text=navy,anchor=south west,inner sep=2pt},
  arr/.style={-{Stealth[length=4pt]},color=slate,thick}]
  \foreach \name/\x/\y/\sty/\txt in {
      corpus/0/0/m/{\fpath{ts/corpus}\\corpus builder}, bench/3.1/0/m/{\fpath{ts/bench}\\test-set builders},
      repo/6.2/0/m/{\fpath{ts/repo.py}\\bridge to the repo}, third/9.3/0/r/{\fpath{third_party/}\\Table-Specialist},
      loop/0/-1.9/m/{\fpath{ts/loop.py}\\TS and NV loop}, repair/3.1/-1.9/m/{\fpath{ts/repair*.py}\\repair-TS},
      ed/6.2/-1.9/m/{\fpath{ts/ed_loop.py}\\\fpath{ts/repair_ed.py}}, rloop/9.3/-1.9/m/{\fpath{ts/r_loop.py}\\NL-to-R},
      gen/0/-3.8/m/{\fpath{ts/gen}\\prompts, extraction}, val/3.1/-3.8/m/{\fpath{ts/validate}\\Algorithm 3},
      train/6.2/-3.8/m/{\fpath{ts/train}\\LoRA SFT, DPO}, llm/9.3/-3.8/m/{\fpath{ts/llm}\\vLLM client, cache},
      eval/0/-5.7/m/{\fpath{ts/eval}\\scorers}, analysis/3.1/-5.7/m/{\fpath{ts/analysis}\\tables, bootstrap},
      sandbox/6.2/-5.7/r/{Docker sandboxes\\Spark, Python, R}, track/9.3/-5.7/m/{\fpath{ts/tracking.py}\\MLflow, run IDs}}
    {\node[\sty] (\name) at (\x,\y) {\txt};}
  \begin{scope}[on background layer]
    \node[lane,fit=(corpus)(third)] (l1) {}; \node[lab] at (l1.north west) {Data and repository};
    \node[lane,fit=(loop)(rloop)] (l2) {};   \node[lab] at (l2.north west) {Experiment drivers};
    \node[lane,fit=(gen)(llm)] (l3) {};      \node[lab] at (l3.north west) {Shared services};
    \node[lane,fit=(eval)(track)] (l4) {};   \node[lab] at (l4.north west) {Evaluation and records};
  \end{scope}
  \draw[arr] (repo) -- (third);
  \draw[arr] (loop) -- (gen); \draw[arr] (repair) -- (val); \draw[arr] (ed) -- (train); \draw[arr] (rloop) -- (llm);
  \draw[arr] (gen) -- (eval); \draw[arr] (eval) -- (analysis); \draw[arr] (val.south) -- (sandbox.north west);
\end{tikzpicture}
\end{adjustbox}
\caption{Module structure, with the main calls shown as arrows. Amber marks repository code, used either through
imports or as patched source; the generators and scorers reach it through \fpath{ts/repo.py}. Every long job runs in its own detached \texttt{tmux} session and takes GPUs and the
validation slot through file locks.}
\label{fig:architecture}
\end{figure}

\section{Provenance: what is inherited and what is ours}\label{sec:provenance}

{\small
\rowcolors{2}{rowshade}{white}
\begin{longtable}{@{}p{3.6cm}p{4.0cm}p{3.6cm}p{3.4cm}@{}}
\caption{Repository components used by our pipeline, how they enter it, and what was changed.}\label{tab:prov-repo}\\
\toprule\rowcolor{headshade}Component & Repository path & How used & Modification\\\midrule\endfirsthead
\toprule\rowcolor{headshade}Component & Repository path & How used & Modification\\\midrule\endhead
\bottomrule\endfoot
Prompt classes and task descriptors (NL-to-SQL, Scala, R, error detection) &
\fpath{Finetune/data_generator/{nl_to_sql,nl_to_sparkscala,nl_to_r,error_detection}/} & Imported from a path added at
run time, for zero-shot test prompts and training formats & None\\
Generation prompts (questions, SQL, Scala, R, typo, semantic error, error validation) &
\fpath{Generation/NL2DSL/*}, \fpath{Generation/ErrorDetection/*}, \fpath{Validation/ErrorDetection/*} & Prompt
strings extracted from each script's syntax tree; scripts never executed (they run on import and hard-code paths) &
None\\
NL-to-SQL evaluator & \fpath{Evaluete/evaluate/ns_evaluator.py} & Class and function definitions compiled from source
(\fpath{ts/repo.py: load_ns_module}) & Patches in \cref{lst:patches}; read-only SQLite\\
Code sandbox & \fpath{table_specialist_sandbox.py}, \fpath{sandbox/Dockerfile.{python,spark}} & Loaded from source
(\fpath{ts/repo.py: load_sandbox}); images built from the repository's Dockerfiles with pinned base digests &
Open-file limit 256 $\to$ 4,096; nothing else\\
Sampling rules & \fpath{Validation/NL2DSL/NL2SQL_SparkScala/prepare_exec.py}, \fpath{execute.py} & Re-implemented to
the same rules & ours, following the repo\\
Scala, R, error-detection evaluators & \fpath{Evaluete/evaluate/{nss,nr,ed}_evaluator.py} & Rules re-implemented
(extraction, reading, equality, metric) & ours, following the repo\\
R sandbox recipe & \fpath{sandbox/Dockerfile.r} & Not used directly; our \fpath{sandbox/Dockerfile.r} adds
\texttt{dplyr} & new file in our tree\\
\end{longtable}
}

\begin{lstlisting}[language=Python,caption={Evaluator patches applied when compiling the repository's
\texttt{ns\_evaluator.py} (\texttt{ts/repo.py}). Loading fails if any pattern stops matching, so an upstream change
cannot silently change the repo-faithful score.},label={lst:patches}]
# pandas 3 removed applymap and errors="ignore"; object dtype keeps fillna("None") working as before.
".applymap(lambda x: pd.to_numeric(x, errors='ignore')).fillna(\"None\")"
    -> ".map(to_numeric_or_keep).astype(object).fillna(\"None\")"
# Fill in the repo's path placeholders (<PATH-TO-TableSpecialist>/...) with our rebuilt test-set folders:
#   WikiSQL cases and databases, Spider, BIRD, WikiTQ cases and databases.
"[[pd.to_numeric(x, errors='ignore')] for x in y_true_sql.split('|')]"
    -> "[[to_numeric_or_keep(x)] for x in y_true_sql.split('|')]"
# sqlite3 is replaced by a wrapper that opens databases read-only.
\end{lstlisting}

{\small
\rowcolors{2}{rowshade}{white}
\begin{longtable}{@{}p{3.0cm}p{5.6cm}p{6.0cm}@{}}
\caption{Components we wrote. Line counts are for the files listed.}\label{tab:prov-ours}\\
\toprule\rowcolor{headshade}Area & Files & Purpose\\\midrule\endfirsthead
\toprule\rowcolor{headshade}Area & Files & Purpose\\\midrule\endhead
\bottomrule\endfoot
Corpus & \fpath{ts/corpus/build.py} & GitTables filtering, de-duplication, contamination check, manifest\\
Test sets & \fpath{ts/bench/{wikisql,wikitq,spider,bird,dev}.py} & Rebuild test and dev sets in the repository's folder
format; check gold against each dataset's engine\\
Model access & \fpath{ts/llm/client.py} & OpenAI-compatible client for vLLM with retries and a SQLite response cache
(\fpath{runs/llm_cache.sqlite})\\
Generation & \fpath{ts/gen/{questions,prompts,code}.py} & Repository prompts; code extraction for SQL, Scala and R\\
Validation & \fpath{ts/validate/exec_agree.py}, \fpath{ts/validate/r_agree.py} & Algorithm 3 for SQL--Scala and
SQL--R, with the comparator\\
Training & \fpath{ts/train/sft.py}, \fpath{ts/train/dpo.py} & LoRA SFT with resumable checkpoints; DPO (RPO loss)\\
Core loop & \fpath{ts/loop.py} & Algorithm 1 for TS and NV; base and Instruct profiles; GPU pool; resume from files\\
NL-to-SQL scoring & \fpath{ts/eval/run_baseline.py}, \fpath{ts/eval/compare.py} & Repo-faithful and corrected scores,
Wilson intervals\\
Other-task scoring & \fpath{ts/eval/{scala_eval,r_eval,ed_eval}.py} & Repository rules for Scala, R and error
detection; deduplicated execution caches\\
Error-detection loop & \fpath{ts/ed_loop.py} & Algorithm 1 with Algorithm 2\\
NL-to-R loop & \fpath{ts/r_loop.py} & Algorithm 1 with SQL--R agreement; TS and repair-TS\\
Extensions & \fpath{ts/ext2.py}, \fpath{ts/ext2_tune.py}, \fpath{ts/repair.py}, \fpath{ts/repair_ts.py},
\fpath{ts/repair_loop.py}, \fpath{ts/data_control.py}, \fpath{ts/repair_ed.py}, \fpath{ts/eval/self_consistency.py},
\fpath{ts/sc_run.py} & bestof4 and DPO; DPO dev tuning; repair-TS rounds 1--3; same-data control; repair-ED and
controls; self-consistency\\
Analysis & \fpath{ts/analysis/{results_table,significance,scala_table}.py} & Results tables and paired bootstrap tests\\
Tracking and set-up & \fpath{ts/tracking.py}, \fpath{ts/a1_instruct_gate.py}, \fpath{scripts/} & MLflow, run-ID
scheme, downloads, smoke tests, throughput measurement\\
Report & \fpath{report/} & This document; \fpath{report/tools/build_tables.py} generates every results table\\
\end{longtable}
}

\section{Reliability engineering}\label{sec:reliability}

The study ran for fourteen days on one machine. In that time the machine crashed once, stalled for 7.5 hours, and lost
its WSL virtual machine once (\cref{app:incidents}). Each mechanism below was added in response to a concrete failure,
and none lost a finished result.

\begin{description}[style=nextline,leftmargin=1.2em]
\item[Resume from files.] Every stage writes its output file last, so the presence of the file marks the stage as
done (questions, code, validation, examples, adapter, scores, summary). A finished round is never rewritten.
\item[Journalled validation.] Each validated pair is appended to a journal and flushed to disk at once (with
\texttt{fsync}), so a power cut loses only the pairs in flight. A truncated last line is skipped.
\item[Frequent checkpoints.] Each run takes about 12 checkpoints, or one every 5 minutes on large runs. Resuming was
tested with hard kills: training resumed from \texttt{checkpoint-55} and finished; validation kept 45 of 58 saved
pairs and finished the rest.
\item[Response cache.] Every LLM request is cached in SQLite, keyed on the model name, prompt and parameters.
Interrupted generation resumes from the cache, and adapters are named by a weight hash so that keys cannot collide
across arms.
\item[File locks.] GPU slots (\fpath{runs/.gpu{0,1}.lock}) and the single CPU validation slot use \texttt{flock}. A
crash leaves no stale lock. One job runs per GPU, and one 36-worker Spark validation runs at a time.
\item[Stall-aware execution cache.] Spark results that timed out after more than 600~s (the limit is 120~s) can only
come from a machine stall. They are re-run and stored beside the original record, which is never overwritten.
\item[Robust JSONL.] Rows are split only on \texttt{\textbackslash n}, because \texttt{str.splitlines} also breaks on
U+2028, which broke a training file.
\item[Operating rules.] Long jobs run in their own detached \texttt{tmux} session with an append-only log. Checkpoints
are copied aside before any resume, and run directories are never deleted, moved or overwritten.
\end{description}

\section{Security of generated code}

Model-written Scala, Python and R run only inside the repository's Docker sandbox:
\begin{itemize}
\item no network, a read-only file system and a non-root user;
\item resource limits;
\item a fresh container per program.
\end{itemize}
Only \texttt{data.csv} and the code file enter the container, and only one result CSV leaves it. Model-written SQL
runs in SQLite on the host through a read-only connection, with SELECT statements only and a 5~s limit.

\section{Testing}

The test suite (\fpath{tests/}, 23 files) covers:
\begin{itemize}
\item prompt extraction against the repository scripts;
\item the comparator, including a 20,000-row speed test;
\item corpus filters, test-set builders and repo-evaluator loading;
\item resumption logic, with fake models;
\item the repair, repair-ED, Scala, R and error-detection scorers.
\end{itemize}
At the time of writing, 205 tests pass and 1 is skipped. The suite needs no GPU. Real-sandbox checks were run
separately, for example the repository's own sandbox integration tests (5 passed, 1 skipped) and hand-written R
programs scored in the R sandbox.
